\documentclass[a4paper]{article}
\usepackage{amsfonts}
\usepackage{amsmath}
\usepackage{amssymb}
\usepackage{graphicx}     

\newcommand{\cis}{\operatorname{cis}}

\begin{document}

\noindent \large Auteur: Hendrik De Bruyn \\
\noindent \large Studierichting: Informatica\\
\noindent \large Oefeningenreeks 1D nr. 24.2\\

\medskip

\normalsize

\textbf{Opgave: } Los de vergelijking $8z^6 + 7iz^3 + 1 = 0$ op in $\mathbb{C}$. \\

\textbf{Oplossing: } \\

Stel $Z = z^3$ dan is $8Z^2 + 7iZ + 1 = 0$. \\

Hieruit volgt dat $Z_{1,2} = \dfrac{-7i \pm \sqrt{(7i)^2 - 4\cdot8}}{16} = \dfrac{-7i \pm 9i}{16} = \left<\begin{array}{ll}
		+: & \dfrac{i}{8} = \dfrac{1}{8} \cis \dfrac{\pi}{2} \\
		-: & -i = \cis \dfrac{3\pi}{2}
	\end{array}\right.$ \\

$Z_1 = z^3 = \dfrac{1}{8} \cis \dfrac{\pi}{2} \Rightarrow \left\{\begin{array}{l}
	z_1 = \dfrac{1}{2} \cis \dfrac{\pi}{6} = \dfrac{\sqrt{3}+i}{4} \\
	z_2 = \dfrac{1}{2} \cis \dfrac{5\pi}{6} = \dfrac{i-\sqrt{3}}{4} \\
	z_3 = \dfrac{1}{2} \cis \dfrac{9\pi}{6} = \dfrac{1}{2} \cis \dfrac{3\pi}{2} = -\dfrac{i}{2}
\end{array}\right.$

$Z_2 = z^3 = \cis \dfrac{3\pi}{2} \Rightarrow \left\{\begin{array}{l}
	z_1 = \cis \dfrac{3\pi}{6} = \cis \dfrac{\pi}{2} = i \\
	z_2 = \cis \dfrac{7\pi}{6} = \dfrac{-\sqrt{3}-i}{2} \\
	z_3 = \cis \dfrac{11\pi}{6} = \dfrac{\sqrt{3}-i}{2}
\end{array}\right.$ \\

\textbf{Oplossing: } De oplossingen voor $8z^6 + 7iz^3 + 1$ zijn $\left\{\dfrac{\sqrt{3}+i}{4}, \dfrac{i-\sqrt{3}}{4}, -\dfrac{i}{2}, i, \dfrac{-\sqrt{3}-i}{2}, \dfrac{\sqrt{3}-i}{2}\right\}$ \\


\begin{thebibliography}{999}
	%\bibitem{a} Referentie
\end{thebibliography}
\end{document}
